A prática de code review tornou-se uma constante nos processos de desenvolvimento
ágeis. No contexto de sistemas open source, as atividades de code review
acontecem a partir da avaliação de contribuições submetidas por meio de
Pull Requests (PR). Este trabalho tem como objetivo analisar a atividade de
code review desenvolvida em repositórios populares do GitHub, identificando
variáveis que influenciam no merge de um PR.

Investigamos sete questões de pesquisa organizadas em duas dimensões:
\begin{itemize}
    \item \textbf{Dimensão A:} relação entre características do PR e seu status final (merged ou closed);
    \item \textbf{Dimensão B:} relação entre características do PR e o número de revisões realizadas.
\end{itemize}

As características analisadas são: tamanho (arquivos e linhas alteradas),
tempo de análise, descrição (número de caracteres) e interações
(número de participantes).
